\subsection{Whitened posterior geometry}
\begin{theoremEnd}[normal,category=geometry]{proposition}[Strong log-concavity of the marginal posterior]
\CombinedResultCheck{prop:geometry}{proposition}\label{prop:geometry}
Under \cref{thm:finite}, let
$\boldsymbol B\boldsymbol B^T=\boldsymbol V_g$ with $\boldsymbol B$
invertible and put $\boldsymbol g=\boldsymbol m_g+\boldsymbol B\boldsymbol z$.
The induced density $q_r$ of $\boldsymbol z$ is $1$-strongly
log-concave. That is, $\log q_r(\boldsymbol z)+\|\boldsymbol z\|^2/2$
is concave. It has a unique mode and finite moments of every
polynomial order.
\end{theoremEnd}
\begin{proofEnd}
\CombinedProofCheck{prop:geometry}
Fix the Gaussian posterior, selected order, and invertible whitening
factor in \cref{prop:geometry}. Since $\tau_r$ is a maximum of
linear functions, it is convex. Subtracting the affine conditional
mean and dividing by $s>0$ therefore makes
$\alpha_r=(\tau_r-\mu)/s$ convex in the whitened coordinates.
Under the ordinary posterior,
$\eta=(\gamma-\mu(\boldsymbol g))/s$ is standard normal and independent
of $\boldsymbol z$. The joint standard normal density in
$(\boldsymbol z,\eta)$ restricted to the convex epigraph
$\eta\geq\alpha_r(\boldsymbol z)$ is $1$-strongly log-concave and
has positive normalizer by \cref{thm:finite}.
Preservation under marginalization
\citep[Theorem 3.8]{saumard2014logconcavity} gives the same constant
for $q_r$.
To obtain the moment and mode conclusions, use the survival factor's
continuity, positivity, and upper bound of one:
$q_r$ is continuous and bounded above by the standard normal density
divided by the positive event probability. This bound gives all
polynomial moments and convergence of the density to zero at infinity,
so continuity implies attainment of a mode, which strong
log-concavity makes unique.
\end{proofEnd}
